\documentclass[12pt]{article}
\usepackage{amsmath, amsthm, amssymb}
\usepackage[margin=2.5cm]{geometry}
\usepackage{hyperref}

\newtheorem{theorem}{Theorem}
\newtheorem{lemma}{Lemma}

\title{\textbf{The Integral Test --- Formal Proof}}
\author{Series and Integral Transforms, Course 21183\\Holon Institute of Technology}
\date{}

\begin{document}
\maketitle

\begin{theorem}[Integral Test]
Let $f:[1,\infty)\to\mathbb{R}$ be a function that is \textbf{positive}, \textbf{continuous},
and \textbf{decreasing} on $[1,\infty)$. Define $a_n = f(n)$. Then
\[
    \sum_{n=1}^{\infty} a_n \quad \text{converges} \quad \iff \quad \int_1^{\infty} f(x)\,dx \quad \text{converges.}
\]
\end{theorem}

\bigskip
\noindent\textbf{Key Inequality.}
Since $f$ is decreasing, for every integer $n \geq 1$ and every $x \in [n,\, n+1]$:
\[
    f(n+1) \leq f(x) \leq f(n).
\]
Integrating over $[n,\, n+1]$ (interval of length 1):
\[
    f(n+1) \;\leq\; \int_n^{n+1} f(x)\,dx \;\leq\; f(n).
\]
Summing from $n = 1$ to $n = N-1$:
\[
    \sum_{n=1}^{N-1} f(n+1)
    \;\leq\; \int_1^N f(x)\,dx
    \;\leq\; \sum_{n=1}^{N-1} f(n),
\]
which gives us the fundamental double inequality:
\begin{equation}\label{eq:key}
    \sum_{n=2}^{N} a_n
    \;\leq\; \int_1^N f(x)\,dx
    \;\leq\; \sum_{n=1}^{N-1} a_n.
\end{equation}

\bigskip
\begin{proof}
Let $S_N = \sum_{n=1}^N a_n$ denote the $N$-th partial sum of the series.
Note that $\{S_N\}$ is \emph{increasing} since all terms $a_n = f(n) > 0$.

\medskip
\noindent\textbf{($\Leftarrow$) If the integral converges, then the series converges.}

Suppose $\int_1^{\infty} f(x)\,dx = L < \infty$. From the right inequality in \eqref{eq:key}:
\[
    \int_1^N f(x)\,dx \;\leq\; \sum_{n=1}^{N-1} a_n = S_{N-1} \leq S_N.
\]
Using the left inequality in \eqref{eq:key}:
\[
    S_N = a_1 + \sum_{n=2}^{N} a_n
        \;\leq\; a_1 + \int_1^N f(x)\,dx
        \;\leq\; a_1 + L.
\]
So $\{S_N\}$ is increasing and bounded above by $a_1 + L$.
By the Monotone Convergence Theorem, $\{S_N\}$ converges, i.e.\ $\sum_{n=1}^{\infty} a_n$ converges.

\medskip
\noindent\textbf{($\Rightarrow$) If the series converges, then the integral converges.}

Suppose $\sum_{n=1}^{\infty} a_n = S < \infty$.
The function $I(t) = \int_1^t f(x)\,dx$ is increasing in $t$ (since $f > 0$).
From the right inequality in \eqref{eq:key}:
\[
    \int_1^N f(x)\,dx \;\leq\; \sum_{n=1}^{N-1} a_n \;\leq\; S.
\]
So $I(N)$ is increasing and bounded above by $S$.
By the Monotone Convergence Theorem, $\lim_{N\to\infty} I(N)$ exists,
i.e.\ $\int_1^{\infty} f(x)\,dx$ converges.
\end{proof}

\bigskip
\noindent\textbf{Remark.} The two directions are completely symmetric in structure:
in both cases we have a monotone increasing quantity (partial sums or integral),
and we use inequality \eqref{eq:key} to bound it above, then apply the
Monotone Convergence Theorem.

\bigskip
\noindent\textbf{Example.} The $p$-series $\displaystyle\sum_{n=1}^{\infty} \frac{1}{n^p}$
with $f(x) = x^{-p}$ (positive, continuous, decreasing for $p>0$, $x\geq 1$).
\[
    \int_1^{\infty} \frac{dx}{x^p} =
    \begin{cases}
        \dfrac{1}{p-1} & \text{if } p > 1 \quad (\text{converges}), \\[6pt]
        +\infty         & \text{if } p \leq 1 \quad (\text{diverges}).
    \end{cases}
\]
Therefore $\sum 1/n^p$ converges if and only if $p > 1$.

\end{document}
